\documentclass[11pt,a4paper]{article}
\usepackage[T1]{fontenc}
\usepackage{lmodern}
\usepackage[margin=28mm]{geometry}
\usepackage{amsmath,amssymb,amsthm,mathtools}
\usepackage{microtype,booktabs,enumitem}
\usepackage[authoryear,round]{natbib}
\usepackage[hidelinks]{hyperref}
\hypersetup{pdftitle={Extremal Klein Four Counts in Semi-Extraspecial 2-Groups},
 pdfsubject={Research preprint version 0.1.0},pdfauthor={Qi Liu}}
\usepackage[nameinlink,noabbrev]{cleveref}
\usepackage{fancyhdr}
\pagestyle{fancy}
\fancyhf{}
\fancyhead[L]{\small Extremal Klein four counts}
\fancyhead[R]{\small Preprint v0.1.0}
\fancyfoot[C]{\thepage}
\setlength{\headheight}{14pt}
\setlength{\emergencystretch}{2em}
\newtheorem{theorem}{Theorem}[section]
\newtheorem{proposition}[theorem]{Proposition}
\newtheorem{lemma}[theorem]{Lemma}
\newtheorem{corollary}[theorem]{Corollary}
\theoremstyle{definition}
\newtheorem{definition}[theorem]{Definition}
\newtheorem{example}[theorem]{Example}
\theoremstyle{remark}
\newtheorem{remark}[theorem]{Remark}
\numberwithin{equation}{section}
\newcommand{\F}{\mathbb F}
\newcommand{\Tr}{\operatorname{Tr}}
\newcommand{\rad}{\operatorname{rad}}
\newcommand{\GL}{\operatorname{GL}}
\newcommand{\SL}{\operatorname{SL}}
\newcommand{\Hom}{\operatorname{Hom}}
\newcommand{\im}{\operatorname{im}}
\newcommand{\Gauss}{\mathcal G}
\newcommand{\Walsh}{\mathcal W}
\newcommand{\Pm}{\mathcal P}
\newcommand{\vfour}{v_4}
\newcommand{\Qstar}{Q_*}
\newcommand{\Hp}{H_+}
\newcommand{\eps}{\varepsilon}
\newcommand{\id}{\mathrm{id}}
\title{Extremal Klein Four Counts in\\ Semi-Extraspecial \(2\)-Groups}
\author{Qi Liu}
\date{Research preprint, version 0.1.0\\17 October 2026}
\begin{document}
\maketitle
\begin{abstract}
Let \(G\) be a semi-extraspecial \(2\)-group with
\(Z(G)=G'=\Phi(G)\cong C_2^2\) and
\(\dim_{\F_2}G/Z(G)=2n\).
We consider the branch with the largest possible number of elements
satisfying \(g^2=1\), including the identity.
We give sharp bounds for the number of actual Klein four subgroups,
with separate formulas for odd and even half-dimension and a small
exception at \(n=2\).
We describe equality through the radical and sign of a sum of three
dual quadratic forms. For every even \(n\geq4\), the groups attaining the
lower bound constitute one isomorphism type: an eight-dimensional
quadratic core together with hyperbolic \(\F_4\)-planes, all sharing the
same two-dimensional target. The corresponding group operation is a
central product. The proof retains every \(\F_2\)-linear refinement and
rules out negative scalar-sign compensation between the core and the
entire semisimple remainder. Individual planes in an initial splitting
of the remainder need not all be positive.
Two explicit examples delimit the inference from maximality and the
scope of the even-parameter classification.
\end{abstract}
\noindent\textit{Keywords:} semi-extraspecial groups; Klein four subgroups;
quadratic maps; bent functions; central products.
\par\medskip
\noindent{\small Preprint v0.1.0; \href{https://creativecommons.org/licenses/by/4.0/}{CC BY 4.0}.
Paper: \href{https://doi.org/10.5281/zenodo.23450418}{doi:10.5281/zenodo.23450418}.\\
Code and Chinese companion: \href{https://doi.org/10.5281/zenodo.23450447}{doi:10.5281/zenodo.23450447}.}

\section{The problem and the main results}\label{sec:main}
Counting elements whose squares are the identity does not determine how
those elements commute. The number of Klein four subgroups measures a
specific part of this additional structure: it counts pairs of distinct
nonidentity involutions whose product is again an involution. For special
\(2\)-groups, squares and commutators are encoded by a quadratic map and
its polar. Requiring every central quotient of order two to be
extraspecial makes all nonzero scalar polars nondegenerate. In the case
of a center of order four there are exactly three scalar components,
so both the element count and the subgroup count can be expressed by a
small, coupled family of quadratic character sums.

We study the branch on which the element count is largest. This is also
the totally orthogonal branch in the group class below. Our question is
how far the Klein four count can vary on this branch, and which groups
attain its lower endpoint. The restriction fixes a natural extremal
condition while leaving the interaction of the three square forms to
be determined.

All groups are finite. A nonabelian \(2\)-group \(E\) is
\emph{extraspecial} if \(Z(E)=E'=\Phi(E)\) has order \(2\).
A group \(G\) is \emph{semi-extraspecial} if \(G/U\) is extraspecial for
every subgroup \(U\) of index \(2\) in \(Z(G)\).
Throughout this paper we impose
\begin{equation}\label{eq:class}
 Z(G)=G'=\Phi(G)=W\cong C_2^2,\qquad
 V=G/W\cong\F_2^{\,2n},\qquad |G|=2^{2n+2},
\end{equation}
and require \(G\) to be semi-extraspecial.
The parity parameter is \(n\), the \emph{half-dimension} of \(V\).

Write
\[
 I(G)=\{g\in G:g^2=1\},\quad s(G)=|I(G)|,\quad
 \vfour(G)=|\{H\leq G:H\cong C_2\times C_2\}|.
\]
Thus \(s(G)\) includes the identity and \(\vfour(G)\) counts subgroups,
not their conjugacy classes. Put
\begin{equation}\label{eq:Mn}
 h=2^n,\qquad N=h^2,\qquad
 M_n=\frac{h^4+9h^3+36h^2-36h+8}{24}.
\end{equation}

\begin{theorem}[Sharp numerical bounds]\label{thm:bounds}
There is no group in \eqref{eq:class} with \(n=1\).
For each \(n\geq2\) the class is nonempty and
\[
 s(G)\leq 2^{2n}+3\cdot2^n=N+3h,
\]
with equality attainable. On the branch \(s(G)=N+3h\), the sharp bounds
are
\begin{equation}\label{eq:bounds}
\begin{array}{ll}
 n=2:& M_n\leq\vfour(G)\leq M_n+2^{3n-2},\\[2pt]
 n\geq3\text{ odd}:&
 M_n-2^{3n-5}\leq\vfour(G)\leq M_n+2^{3n-5},\\[2pt]
 n\geq4\text{ even}:&
 M_n-2^{3n-4}\leq\vfour(G)\leq M_n+2^{3n-2}.
\end{array}
\end{equation}
Every displayed endpoint is attained for every indicated \(n\).
For \(n=2\), the only values on this branch are \(53\) and \(69\).
No assertion is made that every integer between the endpoints occurs.
\end{theorem}

The square map is the \(\F_2\)-quadratic map
\(Q:V\to W,\ Q(gW)=g^2\), with \(W\) written additively.
Two such maps \(Q:V\to W\) and \(\widetilde Q:\widetilde V\to\widetilde W\)
are \emph{strictly equivalent} if there are invertible
\(\F_2\)-linear maps \(S:V\to\widetilde V\) and \(T:W\to\widetilde W\) such that
\begin{equation}\label{eq:strict}
 \widetilde Q(Sv)=TQ(v)\qquad(v\in V).
\end{equation}
In particular, adding an arbitrary linear map is not part of this
equivalence. Its precise group-theoretic meaning is proved in
\cref{prop:strict}.

Let \(K=\F_4=\F_2[\alpha]\), where \(\alpha^2+\alpha+1=0\).
The trace \(\Tr=\Tr_{K/\F_2}\) is regarded as \(K\)-valued through
\(\F_2\subset K\). Define
\begin{equation}\label{eq:representatives}
 \Qstar(a,b,c,d)=ad+bc+\Tr(ac)+a+\alpha c,\qquad
 \Hp(u,v)=uv.
\end{equation}
Both are \(\F_2\)-quadratic maps with target \(K\).
The use of \(K\)-coordinates does \emph{not} declare \(\Qstar\) to be a
homogeneous quadratic form over \(K\).

\begin{theorem}[Even-parameter lower equality]\label{thm:classification}
Assume \eqref{eq:class}, and let \(n\geq4\) be even.
The following conditions are equivalent:
\begin{enumerate}[label=(\roman*)]
\item \(s(G)=2^{2n}+3\cdot2^n\) and
      \(\vfour(G)=M_n-2^{3n-4}\).
\item The square map of \(G\) is strictly equivalent to
\[
 \Qstar\perp \Hp^{\perp r},\qquad r=(n-4)/2,
\]
where all summands take values in the same \(W=K\).
\item \(G\) is isomorphic to the central product, along that common \(W\),
of the group \(G_*\) associated with \(\Qstar\) and \(r\) copies of the
order-\(64\) Heisenberg group \(H(K)\).
\end{enumerate}
Consequently there is exactly one group isomorphism type for each
fixed even \(n\geq4\) satisfying (i).
This states existence of a standard decomposition and uniqueness of the
whole group type, not uniqueness of an embedded core or of a central
decomposition.
\end{theorem}

The first numerical instances of \cref{thm:bounds} are listed in
\cref{tab:small}.
\begin{table}[htb]
\centering
\caption{Sharp Klein four counts on the maximum-\(s\) branch.}
\label{tab:small}
\begin{tabular}{crrrr}
\toprule
\(n\)&\(|G|\)&\(s_{\max}\)&\(\min\vfour\)&\(\max\vfour\)\\
\midrule
2&64&28&53&69\\
3&256&88&431&463\\
4&1024&304&4371&5651\\
\bottomrule
\end{tabular}
\end{table}
The first higher-dimensional instance of \cref{thm:classification}
has \(n=6\), \(\dim_{\F_2}V=12\), and \(|G|=16384\).
The classification does not include odd \(n\), other values of \(s\),
or the intermediate attainable values of \(\vfour\).

The methods belong to existing frameworks. Quadratic maps model special
\(2\)-groups and their central products \citep{KK12,KK15}; alternating
pencils describe the commutator structure \citep{LM18,PFMR14,MPG13}.
The scalar-dual description of \citep[Theorem 2.1,
pp.~509--510]{AFH24} and the general Maiorana--McFarland dual formula of
\citep[equation (4), p.~313]{CMP18} provide the character-sum tools used
below. The \(D\equiv0\) upper endpoint is a direct consequence of the
partial-difference-set theorem in \citep[Theorem 6]{WFW23}, combined with
an elementary subgroup lift count. We give that substitution in
\cref{sec:context}.

The results proved here assemble the numerical bounds, their precise
equality conditions, and the exhaustive lower-equality classification
in \cref{thm:classification}. The selection of the \(96\) refinements of
the equality polar, their local normalization, and the sign restriction
between the core and the entire remainder are short deductions from
the cited formulas and classical quadratic-form criteria. They are not
new general duality or equivalence methods. The distinction from a
classification of polar pencils is that the full square map, including
its linear terms, must be retained. All the arguments needed for the
stated results are included; no computational enumeration is used as a
proof of exhaustiveness.

\section{Square maps, strict equivalence, and counting}\label{sec:prelim}
An \(\F_2\)-quadratic map \(Q:V\to W\) is a map with \(Q(0)=0\) for
which
\[
 B(x,y)=Q(x+y)+Q(x)+Q(y)
\]
is bilinear. Its polar \(B\) is alternating. Linear terms are allowed
and carry information.

\begin{proposition}[The group model]\label{prop:strict}
For the groups under consideration, \(Q(gW)=g^2\) is well-defined,
its polar is the commutator map, and all three nonzero scalar polars
\(\lambda B\), \(\lambda\in W^*\setminus\{0\}\), are nondegenerate.
Conversely, any \(\F_2\)-quadratic map from a nonzero \(V\) to
\(W=\F_2^2\) with these three
polars nondegenerate is the square map of a group satisfying
\eqref{eq:class} and the semi-extraspecial condition.
Two resulting groups are isomorphic if and only if their square maps
are strictly equivalent in the sense of \eqref{eq:strict}.
\end{proposition}
\begin{proof}
The hypotheses \(G'=W\leq Z(G)\), \(W^2=1\), and \(G/W\) elementary
abelian imply that \(G\) has class \(2\) and exponent \(4\).
Changing a lift by an element of \(W\) changes neither its square nor
a commutator. The square identity in a class-\(2\) group gives the polar
formula. For \(U=\ker\lambda\), the center of \(G/U\) modulo \(W/U\)
is the radical of \(\lambda B\). Thus \(G/U\) is extraspecial exactly
when this radical is zero.

For the converse, choose coordinates and write
\[
 Q(v)=\sum_i a_i v_i+\sum_{i<j}c_{ij}v_iv_j.
\]
The bilinear map
\[
 F(v,u)=\sum_i a_iv_iu_i+\sum_{i<j}c_{ij}v_i u_j
\]
has \(F(v,v)=Q(v)\). Define on \(V\times W\)
\begin{equation}\label{eq:group-law}
 (v,w)(u,t)=(v+u,w+t+F(v,u)).
\end{equation}
Bilinearity proves the cocycle identity and associativity. The square
and commutator are \(Q\) and \(B\). Nondegeneracy of the scalar polars
gives \(\rad B=0\), so the center is exactly \(W\).
The span of the values of \(B\) is \(W\), since otherwise a nonzero
\(\lambda\) would annihilate \(B\). Hence the derived subgroup is \(W\);
the identity \(\Phi(G)=G^2G'\) then gives \(\Phi(G)=W\).
The preceding quotient argument proves the semi-extraspecial condition.

A group isomorphism preserves the characteristic center and squares,
and therefore induces \(S,T\) in \eqref{eq:strict}.
Conversely, a basis of \(V\), lifted to \(G\), together with a basis of
the elementary abelian center presents \(G\) by centrality, its
generator squares \(Q(e_i)\), and its commutators \(B(e_i,e_j)\).
Collection puts every element in a form
\(w g_1^{a_1}\cdots g_{2n}^{a_{2n}}\), \(a_i\in\{0,1\}\).
The presentation has at most \(4\cdot2^{2n}\) elements and surjects onto
\(G\), so these are all its relations. Sending the center by \(T\) and
each \(g_i\) to a lift of \(Se_i\) respects these relations by
\eqref{eq:strict}. It gives a surjection between groups of equal order
and hence an isomorphism.
\end{proof}

\begin{lemma}[A common target and a central product]\label{lem:central}
If \(V=V_0\oplus V_1\) and \(B(V_0,V_1)=0\), then
\(Q=Q|_{V_0}\perp Q|_{V_1}\) as maps to the \emph{same} \(W\).
When the restrictions have nondegenerate scalar polars, the corresponding
group is their central product along \(W\).
For \(r+1\) factors this means
\[
 (G_0\times\cdots\times G_r)/\mathcal N,\qquad
 \mathcal N=\{(w_0,\ldots,w_r)\in W^{r+1}:\ \textstyle\sum_iw_i=0\}.
\]
\end{lemma}
\begin{proof}
The assertion about \(Q\) is exactly its polar identity.
Inside the group associated with \(Q\), the preimages of \(V_0,V_1\)
commute, intersect in \(W\), and generate the group. Multiplication from
their external direct product has the stated central kernel. The same
argument applies to any number of summands. This also proves the
square-map version of the standard correspondence
\citep[Lemma 2.4]{KK12}.
\end{proof}

Let \(S=Q^{-1}(0)\), \(z=|S|\), and set
\[
 T_0=|\{(x,y)\in S^2:x+y\in S\}|,\qquad
 J_2=|\{(g,h)\in I(G)^2:gh=hg\}|.
\]
\begin{lemma}[Counting actual Klein four subgroups]\label{lem:count}
One has
\begin{equation}\label{eq:count}
 s=4z,\qquad J_2=16T_0=3s-2+6\vfour(G).
\end{equation}
If \(L\) is the number of two-dimensional \(\F_2\)-subspaces on which
\(Q\) vanishes, then
\begin{equation}\label{eq:zero-planes}
 T_0=3z-2+6L,\qquad \vfour(G)=16L+6z-5.
\end{equation}
\end{lemma}
\begin{proof}
Each zero vector has four involutory lifts, independently for each
member of an ordered pair. For \(x,y\in S\), the lifts commute exactly
when \(B(x,y)=Q(x+y)=0\). This proves the first two equalities.
Among commuting ordered pairs of involutory elements, there are \(s\)
diagonal pairs, \(s-1\) pairs with only the first entry the identity,
and \(s-1\) with only the second. Every remaining pair generates a unique
Klein four group; each such subgroup has six ordered generating pairs.
The identical argument in \(S\), with zero in place of the identity,
gives \(T_0=3z-2+6L\) and then \eqref{eq:zero-planes}.
\end{proof}

\section{A third Walsh moment and the sharp bounds}\label{sec:bounds}
Choose a basis of \(W^*\), write the scalar components as \(q_1,q_2\),
and put \(q_3=q_1+q_2\).
For \(\xi\in V^*\), define
\[
 \Walsh_i(\xi)=\sum_{x\in V}(-1)^{q_i(x)+\xi(x)},\qquad
 a_0=\sum_{i=1}^3\Walsh_i(0),\qquad
 \Pm=\sum_{\xi\in V^*}\prod_{i=1}^3\Walsh_i(\xi).
\]
We use the Fourier convention
\(\widehat f(\xi)=\sum_x f(x)(-1)^{\xi(x)}\).

\begin{lemma}[Quadratic character sums]\label{lem:gauss}
A quadratic function \(d\) on a finite \(\F_2\)-space with \(d(0)=0\)
and polar radical of dimension \(\rho\) has character sum zero if its
restriction to that radical is nonzero. Otherwise, if the ambient
dimension is \(2n\), its character sum is
\(\pm2^{n+\rho/2}\).
For a nondegenerate scalar quadratic form \(q\), every Walsh value has
absolute value \(2^n\).
\end{lemma}
\begin{proof}
On the radical, \(d\) is linear. Translating by a radical vector with
value \(1\) pairs terms with opposite signs. If the restriction is zero,
\(d\) descends to a nondegenerate quadratic form on the quotient, with
\(2^\rho\) preimages for each vector.
A symplectic basis decomposes a nondegenerate alternating polar into
two-dimensional planes. On such a plane the quadratic form is
\(xy+ax+by\), whose character sum is \(2(-1)^{ab}\).
Multiplication over the planes proves the formula.
Adding a linear form to a nondegenerate quadratic form changes no
polar, which proves the Walsh assertion.
\end{proof}

\begin{proposition}[The counting identity]\label{prop:moment}
The maximum possible value of \(s\) is at most \(N+3h\).
It is reached if and only if all three original Gauss sums
\(\Walsh_i(0)\) equal \(h\).
Without assuming this equality, one has
\begin{equation}\label{eq:moment}
 T_0=\frac{N^2+3Na_0+3a_0^2+21N}{64}
          +\frac{3\Pm}{32N}.
\end{equation}
On the branch \(s=N+3h\),
\begin{equation}\label{eq:moment-max}
 \vfour(G)=M_n+\frac{\Pm}{4N}.
\end{equation}
\end{proposition}
\begin{proof}
The indicator identity
\[
 1_S=\frac{1+(-1)^{q_1}+(-1)^{q_2}+(-1)^{q_3}}4
\]
gives \(z=(N+a_0)/4\) and \(s=N+a_0\).
By \cref{lem:gauss}, each original Gauss sum is \(\pm h\), proving the
maximum assertion and its equality criterion.

Put \(U(\xi)=\sum_i\Walsh_i(\xi)\). Then
\(\widehat{1_S}(\xi)=(N\mathbf1_{\xi=0}+U(\xi))/4\).
Fourier orthogonality gives \(T_0=N^{-1}\sum_\xi\widehat{1_S}(\xi)^3\).
Since \(\Walsh_i(\xi)^2=N\),
\[
 U(\xi)^3=7NU(\xi)+6\Walsh_1(\xi)\Walsh_2(\xi)\Walsh_3(\xi).
\]
Moreover \(\sum_\xi U(\xi)=3N\), because each \(q_i(0)=0\).
Thus the numerator after expansion is
\(N^3+3N^2a_0+3Na_0^2+21N^2+6\Pm\), divided by \(64N\).
This proves \eqref{eq:moment}. Substituting \(a_0=3h\) and
\eqref{eq:count}, using \(N=h^2\), gives \eqref{eq:moment-max}.
\end{proof}

For a nondegenerate binary quadratic form, its actual Walsh dual records
the hyperbolic or elliptic type of every linear refinement; this is
\citep[Theorem 2.1, pp.~509--510]{AFH24}. In binary coordinates, a
diagonal square represents a linear function, so the statement includes
the refinements used here. Let \(B_i\) be the polar matrices of \(q_i\)
in the chosen binary coordinates. More explicitly, if
\(\Walsh_i(0)=h(-1)^{\epsilon_i}\) and covectors are columns, translation
by \(B_i^{-1}\xi\) gives
\[
 \Walsh_i(\xi)=h(-1)^{q_i(B_i^{-1}\xi)+\epsilon_i}.
\]
The origin phase \(\epsilon_i\) is part of the actual dual. We now assume
\(s=N+3h\), so all three \(\epsilon_i\) vanish. Completing the square gives
\begin{equation}\label{eq:duals}
 \Walsh_i(\xi)=h(-1)^{q_i^*(\xi)},\qquad
 q_i^*(\xi)=q_i(B_i^{-1}\xi),\qquad q_i^*(0)=0.
\end{equation}
Indeed translation by \(B_i^{-1}\xi\) removes the Fourier linear
term. Define
\begin{equation}\label{eq:DC}
 D=q_1^*+q_2^*+q_3^*,\qquad
 C=B_1^{-1}+B_2^{-1}+(B_1+B_2)^{-1},\qquad
 \rho=\dim\ker C.
\end{equation}
Here \(D\) is the pointwise sum in \(\F_2\), whereas
\(\Gauss(D)=\sum_\xi(-1)^{D(\xi)}\) is an integer.
Thus \(D\equiv0\) gives \(\Gauss(D)=N\), not zero.
The matrix \(C\) is the polar of \(D\). In particular
\begin{equation}\label{eq:P-dual}
 \Pm=h^3\sum_{\xi}(-1)^{D(\xi)}.
\end{equation}
Each inverse \(B_i^{-1}\) is alternating: for \(\xi=B_i x\),
\(\xi^{\mathsf T}B_i^{-1}\xi=x^{\mathsf T}B_ix=0\).

\begin{lemma}[Paired cyclic summands]\label{lem:paired}
Let \(\beta\) be a nondegenerate alternating form over \(\F_2\), and
let \(A\) be self-adjoint for \(\beta\), with \(\beta(v,Av)=0\) for
every \(v\). There is a \(\beta\)-orthogonal decomposition into
\(A\)-invariant nondegenerate spaces \(X\oplus Y\), where \(X,Y\) are
totally isotropic cyclic modules with the same elementary divisor
\(f^k\). In particular every elementary divisor occurs in pairs.
\end{lemma}
\begin{proof}
Self-adjointness gives
\[
\begin{aligned}
 \beta(v,A^{2j}v)&=\beta(A^jv,A^jv)=0,\\
 \beta(v,A^{2j+1}v)&=\beta(A^jv,A^{j+1}v)=0.
\end{aligned}
\]
Every cyclic module is therefore totally isotropic.
Distinct primary components are orthogonal: on one of them a
polynomial annihilating the other is invertible, and self-adjointness
moves that polynomial to the other argument. Hence each primary
component is nondegenerate.

In an \(f\)-primary component choose \(v\) with maximal annihilator
\(f^k\) and let \(X=\F_2[A]v\).
Choose \(w\) in that component with
\(\beta(f(A)^{k-1}v,w)\ne0\), and put \(Y=\F_2[A]w\).
The kernel of the pairing map \(X\to Y^*\) is a submodule.
Any nonzero submodule of \(\F_2[t]/(f^k)\) contains its socle
\(f^{k-1}X\), but the chosen pairing is not zero on that socle.
The kernel is zero. Maximality of \(k\) gives \(\dim Y\leq\dim X\);
thus the pairing is perfect and dimensions are equal.
Total isotropy and perfect pairing imply \(X\cap Y=0\) and
nondegeneracy of \(X\oplus Y\). Its orthogonal complement is
\(A\)-invariant. Repeating proves the lemma.
\end{proof}

\begin{proposition}[Radical restrictions]\label{prop:radical}
Set \(A=B_1^{-1}B_2\) and \(p(t)=t^2+t+1\). Then \(A,I+A\) are
invertible and
\begin{equation}\label{eq:kernel}
 C=[A(I+A)]^{-1}p(A)B_1^{-1},\qquad
 \ker C=B_1\ker p(A).
\end{equation}
Consequently \(\rho\) is divisible by \(4\).
If \(n\) is odd, then \(\rho\leq2n-6\).
If \(n\) is even and \(\sum_\xi(-1)^{D(\xi)}<0\), then
\(\rho\leq2n-4\).
\end{proposition}
\begin{proof}
Writing \(B_2=B_1A\) and \(B_3=B_1(I+A)\) proves invertibility and
\eqref{eq:kernel} by direct multiplication.
Notice that the factor \(B_1\) in the kernel equality must be retained.
For \(\beta(x,y)=x^{\mathsf T}B_1y\), the matrix \(A\) is self-adjoint,
and \(\beta(v,Av)=0\), because \(B_2\) is alternating.
Apply \cref{lem:paired}.
A pair of \(p^k\)-blocks occupies \(4k\) dimensions and contributes
\(4\) to \(\dim\ker p(A)\). A pair belonging to any other irreducible
polynomial contributes zero to this kernel.
There are no linear factors, and \(p\) is the only irreducible
quadratic over \(\F_2\); every other pair therefore occupies at least
six dimensions.

If \(n\) is odd, the identity
\(n=\sum_{\text{pairs}}k\deg f\) forces at least one pair with both
\(k\) and \(\deg f\) odd. Its degree is at least three, so the
non-\(p\) part occupies at least six dimensions.
This gives \(\rho\leq2n-6\).
For even \(n\), if \(\rho=2n\), then \(C=0\) and \(D\) is linear
with \(D(0)=0\). Its sum is \(N\) or \(0\), never negative.
Divisibility by four then yields \(\rho\leq2n-4\) in the negative case.
\end{proof}

\begin{proposition}[Bounds and exact equality tests]\label{prop:equality}
On the maximum-\(s\) branch, the inequalities in \eqref{eq:bounds}
hold. Their equality tests are:
\begin{enumerate}[label=(\alph*)]
\item For odd \(n\geq3\), either endpoint requires and is implied by:
one pair of simple irreducible cubic blocks and all remaining pairs
simple \(p\)-blocks; \(D|_{\ker C}=0\); and the Gauss sum of the
six-dimensional quotient is respectively \(-8\) or \(+8\).
\item For even \(n\geq4\), lower equality requires and is implied by:
one pair of \(p^2\)-blocks and all remaining pairs simple \(p\)-blocks;
\(D|_{\ker C}=0\); and the Gauss sum of the four-dimensional quotient
is \(-4\).
\item For even \(n\geq2\), upper equality holds if and only if \(D=0\).
Equivalently \(p(A)=0\) and the remaining linear function \(D\)
vanishes identically.
\item For \(n=2\), \(C=0\); the lower value \(53\) holds if and only if
\(D\) is a nonzero linear function, and the upper value \(69\) if and
only if \(D=0\).
\end{enumerate}
\end{proposition}
\begin{proof}
By \cref{lem:gauss}, the sum for \(D\) is zero or
\(\pm2^{n+\rho/2}\), with the nonzero case exactly when \(D\) vanishes
on its polar radical. Combining \cref{prop:radical} with
\eqref{eq:P-dual} gives \(|\Pm|\leq2^{5n-3}\) for odd \(n\).
For even \(n\), the positive bound is \(\Pm\leq2^{5n}\);
a negative value satisfies \(\Pm\geq-2^{5n-2}\).
Equation \eqref{eq:moment-max} gives the inequalities.

For the odd endpoints, equality in magnitude forces \(\rho=2n-6\).
The deficit \(2n-\rho\) contributed by a \(p^k\)-pair is \(4(k-1)\);
a non-\(p\) pair contributes \(2k\deg f\geq6\).
The odd-parity argument forces such a non-\(p\) pair, and deficit six
therefore means precisely one simple cubic pair and otherwise simple
\(p\)-pairs. The remaining equality conditions follow from the
quadratic sum on the quotient. Conversely these conditions give exactly
the asserted magnitude and sign.
Deficit four similarly forces exactly one \(p^2\)-pair, proving (b)
in both directions. To attain the even upper bound, every term
\((-1)^{D(\xi)}\) must be \(+1\), so \(D=0\); this also forces
\(p(A)=0\) by \eqref{eq:kernel}.
In dimension four there is only one simple \(p\)-pair, hence \(C=0\);
the two possible linear sums give (d).
In dimension two the space of alternating forms has dimension one,
so three nonzero scalar polars cannot all be nondegenerate.
\end{proof}

\section{Attainment of all numerical endpoints}\label{sec:attain}
The following constructions prove the sharpness assertions without
using the classification theorem. The trace-product positive blocks
are standard Maiorana--McFarland constructions
\citep[Theorem 2(i), p.~314]{CMP18}; the calculations retain the parameter
correspondence and the signs needed here.

\begin{lemma}[The four-dimensional blocks]\label{lem:H}
Identify \(W\) with \(K=\F_4\) by
\begin{equation}\label{eq:iota}
 \iota(w)=(\Tr(w),\Tr(\alpha w)).
\end{equation}
The map \(\Hp(u,v)=uv\) has all scalar Gauss sums \(+4\), dual function \(D\equiv0\),
and \(\Pm=2^{10}\). The map \(uv+u\) has the same polar and the same
\(s=28\), but its Klein four count is \(53\), whereas that of \(\Hp\)
is \(69\).
\end{lemma}
\begin{proof}
For \(z\ne0\),
\[
 \sum_{u,v\in K}(-1)^{\Tr(zuv+\rho u+\sigma v)}
       =4(-1)^{\Tr(\rho\sigma/z)}.
\]
The sum of \(z^{-1}\) over \(K^\times\) is zero, so the dual function
vanishes identically and its integer character sum is \(16\).
For \(uv+u=u(v+1)\), the zero set consists of the four vectors
\((0,v)\) and the three vectors \((u,1)\) with \(u\ne0\).
Its only totally zero two-dimensional \(\F_2\)-subspace is the line
\(u=0\): two distinct nonzero \(u\)'s in the second set have
a sum outside the zero set, and adding a nonzero vector from the
line \(u=0\) also leaves the zero set.
For \(uv\), the zero set is the union of the two coordinate lines,
which are its only totally zero two-dimensional subspaces.
In each case \(z=7\), and \eqref{eq:zero-planes} gives \(53\) and \(69\).
\end{proof}

\begin{lemma}[The six-dimensional endpoint blocks]\label{lem:F8}
Let \(L=\F_8=\F_2[t]/(t^3+t+1)\), with absolute trace \(\Tr_8\).
On \(L^2\) put
\[
 q_1(x,y)=\Tr_8(xy),\qquad
 q_2(x,y)=\Tr_8(txy+cx+dy).
\]
For \(c=d=0\) and for \(c=1,\ d=t^2+t+1\), the three scalar polars are
nondegenerate and all three original Gauss sums are \(+8\).
The respective dual sums have character sums \(+8\) and \(-8\);
therefore \(\Pm=+2^{12}\) and \(-2^{12}\).
\end{lemma}
\begin{proof}
The three polar coefficients are \(1,t,1+t\), all nonzero.
For \(a\ne0\), summing first over \(y\) gives
\begin{equation}\label{eq:field-sum}
 \sum_{x,y\in L}(-1)^{\Tr_8(axy+cx+dy+ux+vy)}
 =8(-1)^{\Tr_8((u+c)(v+d)/a)}.
\end{equation}
When \(c=d=0\), the original sums are \(+8\).
Let \(k=t^{-1}+(1+t)^{-1}=t+1\). The dual sum is then
\(\Tr_8(tuv)\), of character sum \(+8\).
For the second choice, \(d=t^5\),
\(d/t=t^4=t^2+t\), and \(d/(1+t)=t^2\);
both traces are zero, so the original signs again are positive.
The complete dual sum is
\[
 D(u,v)=\Tr_8(tuv+kdu+kcv+kcd).
\]
Using \eqref{eq:field-sum}, its sum is
\[
 8(-1)^{\Tr_8(cd(k+k^2/t))}
 =8(-1)^{\Tr_8(cd\,k/t)}=-8,
\]
because \(k+k^2/t=k/t\), \((k/t)d=1\), and \(\Tr_8(1)=1\).
This also shows that the polar of this \(D\) is nondegenerate.
\end{proof}

\begin{lemma}[An eight-dimensional lower block]\label{lem:attain8}
There exists a square map on \(\F_2^8\) with all original scalar Gauss
sums \(+16\) and with dual character sum \(-64\).
\end{lemma}
\begin{proof}
Let \(U\) be multiplication by \(t\) on
\(\F_2[t]/((t^2+t+1)^2)\), in its polynomial basis, and put
\[
 E=p(U),\qquad K_0=U^{-1}+(I+U)^{-1}.
\]
Here \(E^2=0\), \(\operatorname{rank}E=2\), and
\(I+K_0=E\), since \(U(I+U)=I+E\).
Indeed multiplication by \(p\) on \(\F_2[t]/(p^2)\) has image and
kernel equal to the ideal \((p)\), each of binary dimension two.
Choose \(0\ne c\in\im E\), \(z\) with \(z^{\mathsf T}c=1\), and
\(d=E^{\mathsf T}z\). On \(\F_2^4\oplus\F_2^4\), take
\[
 q_1(x,y)=x^{\mathsf T}y,\qquad
 q_2(x,y)=x^{\mathsf T}Uy+c^{\mathsf T}x+d^{\mathsf T}y.
\]
The matrices \(I,U,I+U\) are invertible, proving nondegeneracy of the
three scalar polars. The two possibly nontrivial original Gauss signs
are \(d^{\mathsf T}U^{-1}c\) and
\(d^{\mathsf T}(I+U)^{-1}c\); both vanish because \(Ec=0\) and \(E\)
commutes with these matrices. Thus all sums are \(+16\).

The same summation used in \eqref{eq:field-sum}, now with ordinary
binary dot products, gives the complete dual sum at covectors
\((a,b)\) as
\[
 D(a,b)=b^{\mathsf T}(Ea+c)+z^{\mathsf T}Ea.
\]
Here \(K_0c=c\), \(EK_0=E\), and \(Ec=0\) account for the linear and
constant terms. Summing over \(b\) leaves precisely the four solutions
of \(Ea=c\). For each, \(z^{\mathsf T}Ea=1\).
The sum is therefore \(-16\cdot4=-64\).
\end{proof}

\begin{proof}[Completion of the proof of \cref{thm:bounds}]
For a common-target orthogonal sum, each scalar Walsh transform is a
product of the corresponding transforms. The third moment \(\Pm\)
and each original Gauss sum consequently multiply.
For even \(n\), \(n/2\) copies of \(\Hp\) have original sums \(h\)
and \(\Pm=2^{5n}\), giving the upper endpoint. Equivalently, in the notation of
\citep[Theorem 2(i)]{CMP18}, take \(p=2\), \(s=2\), \(m=n\),
\(E=\F_{2^n}\), and \(\pi(y)=y\), with \(g=0\).
Then \(\Tr_{E/K}(xy)\) becomes \(\Hp^{\perp n/2}\) in a pair of
trace-dual \(K\)-bases. The condition \(2\mid n\) is essential to this
subfield description.
For odd \(n\), one of the blocks in \cref{lem:F8} together with
\((n-3)/2\) copies of \(\Hp\), all with the same ordered target,
has \(\Pm=\pm2^{5n-3}\).
For even \(n\geq4\), the block in \cref{lem:attain8} with
\((n-4)/2\) copies of \(\Hp\) has \(\Pm=-2^{5n-2}\).
All original Gauss sums in these constructions are \(+h\).
Every scalar polar remains nondegenerate, so \cref{prop:strict}
realizes actual groups in the asserted class.
Finally \cref{lem:H} supplies both \(n=2\) endpoints.
Together with \cref{prop:equality}, this proves every assertion of
\cref{thm:bounds}.
\end{proof}
\section{The eight-dimensional square maps}\label{sec:core}
In this section \(p(t)=t^2+t+1\). We first normalize the ordered
polar pair, keeping the target coordinates fixed.

\begin{lemma}[A simultaneous polar model]\label{lem:polar-model}
Suppose \(A=B_1^{-1}B_2\) has one pair of \(p^2\)-blocks and \(r\)
pairs of \(p\)-blocks. There is a common orthogonal decomposition
\[
 V=V_0\oplus V_1,\qquad \dim_{\F_2}V_0=8,\quad
 \dim_{\F_2}V_1=4r,\quad B(V_0,V_1)=0.
\]
Put \(R=\F_2[t]/(p^2)\), \(\eps=p(t)\), and
\(\alpha=t+\eps=t^2+1\). Then
\(R=K[\eps]/(\eps^2)\), \(K=\F_4\).
After input changes and the single common target identification
\eqref{eq:iota}, the restrictions of \(Q\) have the forms
\begin{align}
 Q^{(0)}={}&ad+bc+\Tr(ac)+fa+ga^2+\ell c+mc^2
                    +\eta b+kb^2+\theta d+jd^2,\label{eq:all-core}\\
 Q^{(1)}={}&\sum_{\nu=1}^r
   (u_\nu v_\nu+F_\nu u_\nu+G_\nu u_\nu^2
                     +L_\nu v_\nu+M_\nu v_\nu^2),\label{eq:all-ss}
\end{align}
where all coefficients lie in \(K\).
These include all \(\F_2\)-linear refinements of the fixed polars.
Each nonzero scalar polar of each nonempty restriction is
nondegenerate.
\end{lemma}
\begin{proof}
Use \cref{lem:paired} to select the paired \(p^2\)-cyclic modules
as \(V_0\). Their \(\beta=B_1\)-orthogonal complement \(V_1\) is
\(A\)-invariant by self-adjointness. Hence the decomposition is
orthogonal for both \(B_1\) and \(B_2\), and thus for \(B\).
Each scalar polar is a nondegenerate block diagonal form.
The polar identity already gives \(Q=Q^{(0)}\perp Q^{(1)}\);
it does not yet imply any sign condition on the restrictions.

Let \(\lambda:R\to\F_2\) extract the coefficient of \(t^3\).
The pairing \(\lambda(xy)\) is nondegenerate: if a polynomial
\(x\ne0\) has degree \(j\leq3\), multiplying by \(t^{3-j}\) gives
coefficient \(1\) at degree three without a reduction.
On a paired cyclic summand \(R^2\), the balanced pairing for \(B_1\)
has the form \(\lambda'(xy'+x'y)\).
Nondegeneracy writes \(\lambda'(z)=\lambda(ez)\) for a unit \(e\in R\);
otherwise multiplication by \(e\) would have a kernel.
Rescaling one input generator by \(e^{-1}\) normalizes the pair to
\begin{equation}\label{eq:R-polar}
 B((x,y),(x',y'))=
 \bigl(\lambda(xy'+x'y),\lambda(t(xy'+x'y))\bigr).
\end{equation}
This uses only input changes. On each simple \(p\)-pair the same
argument, using the nondegenerate trace pairing on \(K\), gives
\((\Tr(uv'+u'v),\Tr(\alpha(uv'+u'v)))\), in the same ordered target.

The relations \(\eps^2=0\), \(\alpha^2+\alpha+1=0\) give the asserted
description of \(R\). Direct evaluation on the basis
\(1,\alpha,\eps,\alpha\eps\) gives
\[
 \lambda(a+\eps b)=\Tr(b),\qquad
 \lambda(t(a+\eps b))=\Tr(a+\alpha b).
\]
For \(x=a+\eps b,\ y=c+\eps d\), the base refinement
\((\lambda(xy),\lambda(txy))\) therefore becomes
\(ad+bc+\Tr(ac)\) under \(\iota^{-1}\).
On a simple pair the base refinement becomes \(uv\).

Two \(\F_2\)-quadratic maps with the same polar differ by an
\(\F_2\)-linear map. Every such map \(K\to K\) is uniquely
\(x\mapsto ux+vx^2\): the map from the two coefficients is injective
by evaluating at \(1,\alpha\), and both spaces have \(\F_2\)-dimension
four. Applying this independently to all coordinates proves
\eqref{eq:all-core} and \eqref{eq:all-ss}.
No assumption of \(K\)-homogeneity has been made.
\end{proof}

The next formula is the specialization of the general
Maiorana--McFarland formula in \citep[equation (4), p.~313]{CMP18}.
The common \(\F_{16}\) coordinates and the arbitrary function \(g\)
in that substitution are recorded in \cref{app:MM}. We also give a
direct sum calculation to fix the Fourier convention and the constants.

\begin{lemma}[The complete core phases]\label{lem:core-phase}
For \eqref{eq:all-core} and \(z\in K^\times\), let \(\Walsh_z\) be
the Walsh transform of \(\Tr(zQ^{(0)})\). Use the Fourier
functional \(\Tr(\rho a+\sigma b+\tau c+\upsilon d)\), and write
\(\Walsh_z=16(-1)^{E_z}\). With \(\kappa_z=\Tr(z)\), one has
\begin{align}
 a_z&=\upsilon/z+\theta+zj^2,\qquad
 c_z=\sigma/z+\eta+zk^2,\label{eq:az}\\
 E_z&=\Tr\bigl(\kappa_z a_zc_z+
       (\rho+zf+z^2g^2)a_z+(\tau+z\ell+z^2m^2)c_z\bigr).
                                                    \label{eq:raw-core}
\end{align}
Phase additions, including \(\sum E_z\) and \(\gamma\) below, are in
\(\F_2\); sums of \((-1)^{E_z}\) are integer character sums.
In particular the \emph{unadjusted} sum of these three phases is
\begin{align}
 \sum_{z\ne0}E_z
 &=\Tr\bigl(\sigma\upsilon+(f+\eta)\upsilon
          +(\ell+\theta)\sigma+\theta\rho+\eta\tau\bigr)+\gamma,
                                                        \label{eq:core-sum}\\
 \gamma&=\Tr(\theta k^2+\eta j^2+jk+gj+mk)
         =\sum_{z\ne0}E_z(0).                         \label{eq:gamma}
\end{align}
These formulas hold without a hypothesis on the original Gauss signs.
\end{lemma}
\begin{proof}
The identity \(\Tr(za^2)=\Tr(z^2a)\) over \(K\) moves each
squared-variable linear term into the ordinary Fourier coefficient.
Summing first over \(b,d\) forces exactly the two equations in
\eqref{eq:az} and leaves a factor \(4^2=16\).
The remaining exponent is \eqref{eq:raw-core}.
For the sum over \(z\in K^\times\), use
\[
 \sum z=\sum z^2=\sum z^{-1}=0,\quad \sum z^3=1,\quad
 \sum\kappa_z=0,
\]
and
\[
 \sum\kappa_z/z=\sum\kappa_z/z^2
 =\sum\kappa_z z=\sum\kappa_z z^2=1.
\]
Expansion gives the nonconstant terms in \eqref{eq:core-sum}.
The constant terms are
\(\Tr(\theta k^2+\eta j^2+j^2k^2+g^2j^2+m^2k^2)\),
which is \eqref{eq:gamma} since the trace is invariant under squaring.
This proves the formula including its constant.
\end{proof}

\begin{lemma}[The anisotropic sign obstruction]\label{lem:obstruction}
Suppose \(\eta=\theta=0\) in \eqref{eq:all-core}, and all three
original Gauss sums have the same sign \((-1)^\delta\).
Then
\begin{equation}\label{eq:comp}
 fj^2+\ell k^2=jk,\qquad \Tr(gj+mk+jk)=\delta.
\end{equation}
If additionally \(\Tr(f\ell)=1\), then \(j=k=0\) and \(\delta=0\).
\end{lemma}
\begin{proof}
At the Fourier origin, \(a_z=zj^2\) and \(c_z=zk^2\).
Put \(e=jk,\ u=fj^2+\ell k^2,\ v=gj+mk\).
Using
\(\Tr(z)\Tr(ze)=\Tr(z^2e)+\Tr(e)\) for \(z\ne0\), formula
\eqref{eq:raw-core} becomes
\[
 E_z(0)=\Tr(z^2(u+e))+\Tr(v+e).
\]
The three values equal \(\delta\) precisely when \(u+e=0\) and
\(\Tr(v+e)=\delta\): a nonzero linear functional on \(K\) is not
constant on \(K^\times\).
This proves \eqref{eq:comp}.

If \(\Tr(f\ell)=1\), then \(f\ne0\), and the binary \(K\)-quadratic
form \(fX^2+XY+\ell Y^2\) is anisotropic.
Indeed a nonzero solution with \(Y\ne0\) would give, on setting
\(w=fX/Y\), the equation \(w^2+w+f\ell=0\), contradicted by its trace.
If \(Y=0\), then \(X=0\).
The first equation in \eqref{eq:comp} therefore forces \(j=k=0\);
the second now forces \(\delta=0\).
\end{proof}

With the equality polar fixed, the following selection is a short
application of the preceding dual formula, the radical criterion in
\cref{lem:gauss}, and the binary Arf test in \cref{lem:obstruction}.

\begin{proposition}[The \(96\) core refinements]\label{prop:96}
In dimension eight, the maps with maximum \(s=304\) and
\(\vfour=4371\), in the fixed polar coordinates
\eqref{eq:all-core}, are exactly
\begin{equation}\label{eq:96}
 Q_{f,g,\ell,m}=ad+bc+\Tr(ac)+fa+ga^2+\ell c+mc^2,
       \qquad \Tr(f\ell)=1.
\end{equation}
There are \(96\) such coordinate maps.
Every dimension-eight lower-equality map can be put in these
coordinates.
\end{proposition}
\begin{proof}
By \cref{prop:equality}(b), every lower-equality polar has exactly
one pair of \(p^2\)-blocks. Thus \cref{lem:polar-model} covers all of
them, rather than just a selected construction.
Maximum \(s\) means \(E_z(0)=0\) for all \(z\ne0\), hence \(\gamma=0\).
The radical of the polar in \eqref{eq:core-sum} consists of the
directions \(\rho,\tau\), with \(\sigma=\upsilon=0\).
The radical-vanishing condition therefore forces
\(\eta=\theta=0\).
Summing over \(\rho,\tau\) and then completing the square in
\(\sigma,\upsilon\) gives
\begin{equation}\label{eq:core-G}
 \sum_{\rho,\sigma,\tau,\upsilon}(-1)^D
       =64(-1)^{\Tr(f\ell)}.
\end{equation}
Lower equality requires \(-64\), hence \(\Tr(f\ell)=1\).
Apply \cref{lem:obstruction} with \(\delta=0\) to get \(j=k=0\).
This proves necessity of \eqref{eq:96}.
Conversely, for \eqref{eq:96}, formula \eqref{eq:raw-core} at the
origin gives all three signs positive, and \eqref{eq:core-G} is
\(-64\). The counting identity gives
\(\vfour=M_4-2^8=4627-256=4371\).
There are three choices for \(f\ne0\), two for \(f\ell\) with trace
one, and four each for \(g,m\), giving \(3\cdot2\cdot4^2=96\).
\end{proof}

The local input changes below are instances of the transformation
framework of \citep[Propositions 5--6, pp.~12--14]{MPG13} and
\citep[Propositions 4.5 and 4.7, p.~17]{PFMR14}.
The last distinction is the classical binary Arf class over \(K\).
We retain the explicit changes, including their induced linear terms.

\begin{proposition}[Explicit transitivity]\label{prop:orbit}
All \(96\) maps in \eqref{eq:96} are strictly equivalent to \(\Qstar\),
using only input transformations that preserve both scalar polars
and leave the target fixed.
\end{proposition}
\begin{proof}
Every matrix in \(\SL_2(R)\), acting on
\((x,y)=(a+\eps b,c+\eps d)\), preserves
\(xy'+x'y\) and therefore both forms in \eqref{eq:R-polar}.
Starting with \eqref{eq:96}, make the following pullback changes.
\begin{enumerate}[label=\arabic*.]
\item The matrix
\[
 M_1=\begin{pmatrix}1&\eps m\\ \eps g&1\end{pmatrix}
\]
has determinant one and changes \(b\) to \(b+mc\), \(d\) to \(d+ga\).
Substitution cancels both \(ga^2\) and \(mc^2\).
\item The matrix \(M_2=\operatorname{diag}(f^{-1},f)\) changes
the two remaining coefficients to \(1,f\ell\).
The latter is \(\alpha\) or \(\alpha^2\).
\item If \(f\ell=\alpha^2\), use
\[
 M_3=\begin{pmatrix}1&\alpha\\0&1\end{pmatrix}.
\]
The new coefficient of \(c\) is
\(\alpha^2+\alpha+\alpha^2=\alpha\), and an extra
\(\alpha c^2\) is introduced.
The matrix \(M_4=\begin{pmatrix}1&\eps\alpha\\0&1\end{pmatrix}\)
cancels this last term.
\end{enumerate}
The resulting map is exactly \eqref{eq:representatives}.
All the operations act on the entire quadratic map. In particular
the linear correction generated by pulling back the base map
\(ad+bc+\Tr(ac)\) has been included, not discarded.
Each operation has target transformation \(T=\id_W\).
\end{proof}

\begin{remark}\label{rem:isometry}
In fact the full componentwise isometry group of \eqref{eq:R-polar}
is \(\SL_2(R)\). Such an isometry commutes with \(A\), hence is
\(R\)-linear. Its determinant \(d\) satisfies
\(\lambda((d-1)r)=0\) for every \(r\in R\), forcing \(d=1\).
There are \(16^2-4^2=240\) possible unimodular first columns, and
\(16\) second columns with determinant one for each: a unimodular
column has a unit coordinate, so its determinant map \(R^2\to R\)
is onto and each fibre has size \(16\). Thus its order is
\(3840\). This observation is not needed to infer transitivity from
the explicit changes above. Allowing additional target changes cannot
split an already single orbit. Thus no uncomputed subgroup orbit is
being mistaken for a full classification.
\end{remark}

For a concrete group \(G_*\), take \(K^4\times K\) and use
\eqref{eq:group-law} with
\begin{equation}\label{eq:explicit-G}
 F((a,b,c,d),(a',b',c',d'))=
 ad'+bc'+\Tr(ac')+(aa')^2+\alpha(cc')^2.
\end{equation}
This is \(\F_2\)-bilinear, and its diagonal is \(\Qstar\).
The scalar polars in \eqref{eq:R-polar} are nondegenerate because
\(1,t,1+t\) are units in \(R\).
Consequently \cref{prop:strict} and
\cref{prop:96,prop:orbit} prove \cref{thm:classification} for \(n=4\),
including existence of the group of order \(1024\).

\section{The common core in every even half-dimension}
\label{sec:classification}
For the remainder of this section let \(n=4+2r\), \(r\geq0\).
Use the common-target coordinates of \cref{lem:polar-model}.

\begin{lemma}[The semisimple phases]\label{lem:ss-phase}
For \eqref{eq:all-ss}, the scalar Walsh transform at the functional
\(\Tr\sum_\nu(\rho_\nu u_\nu+\sigma_\nu v_\nu)\) is
\(4^r(-1)^{E_{1,z}}\), where
\begin{equation}\label{eq:ss-phase}
 E_{1,z}=
 \sum_\nu\Tr\left(
 \frac{(\rho_\nu+zF_\nu+z^2G_\nu^2)
       (\sigma_\nu+zL_\nu+z^2M_\nu^2)}{z}\right).
\end{equation}
Its phase sum is
\begin{equation}\label{eq:ss-sum}
 D_1=\sum_{z\ne0}E_{1,z}
     =\Tr\sum_\nu(L_\nu\rho_\nu+F_\nu\sigma_\nu+G_\nu M_\nu).
\end{equation}
If \(F_\nu=L_\nu=0\) for all \(\nu\), then each original Gauss sum has
the same sign \((-1)^\delta\), where
\(\delta=\Tr\sum_\nu G_\nu M_\nu\), and \(D_1\equiv\delta\).
\end{lemma}
\begin{proof}
This is also \citep[equation (4), p.~313]{CMP18}: on each plane use
\(p=2\), \(m=2\), \(s=1\), scalar parameter \(1\), the pairing
\(\Tr(\rho u+\sigma v)\), and
\[
 \pi_z(v)=zv+zF+z^2G^2,\qquad
 g_z(v)=(zL+z^2M^2)v.
\]
The permutation \(\pi_z\) is affine; the original squared-variable
terms have been moved by the trace identity, not omitted.
Substitution of \(\pi_z^{-1}(\rho)\) in that formula gives
\eqref{eq:ss-phase}, with its origin phase intact.
Summing its expansion over \(K^\times\), the terms with \(z,z^2,z^{-1}\)
vanish and \(z^3=1\); trace invariance under squaring gives
\eqref{eq:ss-sum}.
With \(F_\nu=L_\nu=0\), evaluating \eqref{eq:ss-phase} at the origin
gives \(\Tr\sum_\nu G_\nu M_\nu=\delta\) for each \(z\).
\end{proof}

The lower-equality sign is essential in the following deduction from
the two phase formulas. Positivity of the whole map alone does not
give positivity on each restriction.

\begin{lemma}[Exclusion of core--remainder negative-sign compensation]\label{lem:no-comp}
Suppose \(Q=Q^{(0)}\perp Q^{(1)}\) is in the form
\eqref{eq:all-core}--\eqref{eq:all-ss}, all three original Gauss sums
are \(+2^n\), and
\begin{equation}\label{eq:lower-dual}
 D|_{\ker C}=0,\qquad
 \sum_{\xi\in V^*}(-1)^{D(\xi)}=-2^{2n-2}.
\end{equation}
Then
\[
 \eta=\theta=j=k=0,\quad \Tr(f\ell)=1,\quad
 F_\nu=L_\nu=0\ (\text{all }\nu),\quad
 \Tr\sum_\nu G_\nu M_\nu=0.
\]
In particular both restrictions separately have all three original
Gauss sums positive.
\end{lemma}
\begin{proof}
Write the \emph{unadjusted} phases on the two blocks as \(E_{0,z}\)
and \(E_{1,z}\). The product rule for Walsh transforms and global
maximality imply only
\begin{equation}\label{eq:matching-signs}
 E_{0,z}(0)=E_{1,z}(0)\qquad(z\ne0).
\end{equation}
The two sides are not yet known to be zero or to be independent of \(z\).
The global normalized dual sum is
\(D=D_0+D_1\), where \(D_j=\sum_{z\ne0}E_{j,z}\).
It has \(D(0)=0\), although its block summands may have constants.

The polar of \(D_1\) is zero by \eqref{eq:ss-sum}, so the whole
space \(V_1^*\) is in the radical of the global polar \(C\).
The equation \(D(0,\xi_1)=0\) for every \(\xi_1\in V_1^*\) makes
\(D_1\) constant. Trace nondegeneracy in \eqref{eq:ss-sum} forces
\(F_\nu=L_\nu=0\) for all \(\nu\).
By \cref{lem:ss-phase}, put
\(\delta=\Tr\sum_\nu G_\nu M_\nu\); then
\(D_1\equiv\delta\) and \(E_{1,z}(0)=\delta\) for every \(z\).
Equation \eqref{eq:matching-signs} now gives
\(E_{0,z}(0)=\delta\), and hence \(\gamma=\delta\) in
\eqref{eq:core-sum}.

The remaining radical directions are the core coordinates
\(\rho,\tau\). On them the global \(D\) is
\(\Tr(\theta\rho+\eta\tau)+\delta+\delta\).
Radical vanishing forces \(\eta=\theta=0\).
The constants cancel between the blocks; they have not been
individually suppressed. The complete phase is now
\[
 D=\Tr(\sigma\upsilon+f\upsilon+\ell\sigma),
\]
independent of \(V_1^*\), \(\rho\), and \(\tau\).
Its character sum is
\[
 2^{4r}\,64(-1)^{\Tr(f\ell)}.
\]
The \emph{negative} sign in \eqref{eq:lower-dual} therefore forces
\(\Tr(f\ell)=1\).
The original core signs all equal \((-1)^\delta\), so
\cref{lem:obstruction} forces \(j=k=0\) and \(\delta=0\).
This proves every assertion. Empty sums give \(\delta=0\) when \(r=0\).
\end{proof}

\begin{lemma}[The positive remainder]\label{lem:hyperbolic}
A map
\begin{equation}\label{eq:positive-remainder}
 q(u,v)=\sum_{\nu=1}^r
 (u_\nu v_\nu+G_\nu u_\nu^2+M_\nu v_\nu^2),\qquad
 \Tr\sum_\nu G_\nu M_\nu=0,
\end{equation}
is isometric over \(K\), with target unchanged, to
\(\Hp^{\perp r}\).
\end{lemma}
\begin{proof}
At this stage, unlike \eqref{eq:all-ss}, the map really is a
homogeneous quadratic form over \(K\).
The hyperbolic/elliptic classification recalled in \citep[p.~509]{AFH24}
now applies over \(K\); equivalently its component signs can be read
from \citep[Theorem 2.2, pp.~511--512]{AFH24} with a trace-compatible
pairing. For \(r\geq1\), identify \(K^{2r}\) linearly with
\(\F_{4^{2r}}\): the theorem's base field is \(K\) and its extension
degree is \(2r\). The required \(K\)-homogeneity has just been established,
rather than inferred from the notation in \cref{lem:polar-model}.
Its polar is the nondegenerate \(K\)-symplectic form
\(\sum_\nu(u_\nu v'_\nu+u'_\nu v_\nu)\).
By \cref{lem:ss-phase}, all three scalar Gauss sums are
\(+4^r\), and its zero count is
\[
 \frac{4^{2r}+3\cdot4^r}{4}.
\]
For \(r\geq1\) this exceeds one, so choose a nonzero isotropic vector
\(u\). Choose \(v\) with \(B_K(u,v)=1\) and replace it by
\(v+q(v)u\). Homogeneity and \(q(u)=0\) show that the new vector is
isotropic and still pairs to one with \(u\). They span a hyperbolic
\(K\)-plane. Its \(K\)-orthogonal complement has nondegenerate polar
and all scalar Gauss sums \(+4^{r-1}\), by dividing out the three
\(+4\) factors. Repetition yields \(r\) hyperbolic planes.
The resulting input change is \(K\)-linear symplectic and fixes \(W\).
The case \(r=0\) is the empty space.
\end{proof}

\begin{proof}[Proof of \cref{thm:classification}]
Assume (i). By \cref{prop:equality}(b), the polar has one pair of
\(p^2\)-blocks and \(r=(n-4)/2\) simple \(p\)-pairs, the global
\(D\) vanishes on its radical, and its character sum is
\(-2^{2n-2}\).
The common model from \cref{lem:polar-model} applies.
By \cref{lem:no-comp}, the core is in \eqref{eq:96} and the remainder
is in \eqref{eq:positive-remainder}.
Normalize the core by \cref{prop:orbit} and the remainder by
\cref{lem:hyperbolic}. Both input changes have \(T=\id_W\).
Their block diagonal sum is thus a legitimate transformation of the
whole square map, proving (ii).

Conversely, the three scalar polars of
\(\Qstar\perp\Hp^{\perp r}\) are blockwise nondegenerate.
Its original Gauss sums are \(16\cdot4^r=2^n\), so \(s=N+3h\).
Its dual character sum is
\(-64\cdot16^r=-2^{2n-2}\).
Thus \(\Pm=-2^{5n-2}\), and \eqref{eq:moment-max} gives exactly
\(\vfour=M_n-2^{3n-4}\). Strict equivalence preserves these quantities.

The group \(H(K)\) can be written on \(K^2\times K\) with
\[
 ((u,v),w)((u',v'),w')=
 ((u+u',v+v'),\,w+w'+uv').
\]
It has square map \(uv\), order \(64\), and center \(K\).
By \cref{lem:central,prop:strict}, (ii) is equivalent to (iii).
In the central product the kernel has order \(4^r\), and the order is
\[
 1024\cdot64^r/4^r=2^{10+4r}=2^{2n+2}.
\]
This proves all directions and the uniqueness of the group type.
It asserts no uniqueness of a chosen core subspace.
\end{proof}

\begin{remark}[Target changes and mixed input changes]\label{rem:target}
The preceding proof already normalized both restrictions with the same
target fixed. Center identifications do not create extra types here:
every member of \(\GL_{\F_2}(K)\) is \(T(w)=c w^\sigma\), with
\(c\in K^\times\) and \(\sigma=\id\) or the Frobenius square.
These are six distinct maps, all the elements of \(\GL_2(\F_2)\).
Each has a lift on a positive plane,
\[
 (u,v)\longmapsto(cu^\sigma,v^\sigma),\qquad
 \Hp(cu^\sigma,v^\sigma)=T(\Hp(u,v)).
\]
Thus a different center gluing on each positive factor can be absorbed
by its input change, without imposing an incompatible target change
on the core.

The remainder in \cref{lem:hyperbolic} need not consist of positive
planes in an initially chosen splitting.
For \(A_-(u,v)=uv+u^2+\alpha v^2\), each scalar original sign is
negative, but \(A_-\perp A_-\) is isometric to
\(\Hp\perp\Hp\) by the same lemma.
For example, a \(K\)-linear change with columns
\[
 (0,1,0,1),\quad(0,1,1,1),\quad(1,0,1,0),\quad(\alpha,1,\alpha,0)
\]
pulls the former back to the latter. Permitting such mixed input
changes is necessary; initial plane signs are not invariants of
the whole remainder. In particular \(\Qstar\perp A_-\perp A_-\)
also attains the corresponding lower bound, by this target-fixed
isometry and \cref{thm:classification}.
\end{remark}

\section{Two distinct limitations of stronger conclusions}\label{sec:examples}
\begin{example}[Maximality does not pass to orthogonal blocks]
\label{ex:compensation}
On \(K^4\) and \(K^2\), with the same target \(K\), take
\[
 Q_0^-=ad+bc+\Tr(ac)+\alpha a^2+d^2,\qquad
 A_-(u,v)=uv+u^2+\alpha v^2.
\]
The three scalar polars of each are nondegenerate.
In the core parameters these are \(g=\alpha,\ j=1\), with
\(f=\ell=m=\eta=k=\theta=0\).
Formula \eqref{eq:raw-core} gives each original core sign negative,
and \cref{lem:ss-phase} gives each original plane sign negative.
Their Gauss sums are therefore respectively \(-16\) and \(-4\),
while the product \(Q_0^-\perp A_-\) has three original sums \(+64\).
The two restrictions have \(s=208\) and \(s=4\), although the whole
group of order \(16384\) has the maximum \(s=4288\).

Here \(D_0=\Tr(\sigma\upsilon)+1\) and \(D_1=1\).
The global dual character sum is \(+1024\), not \(-1024\).
Since \(M_6=803403\), equation \eqref{eq:moment-max} gives
\[
 \vfour(Q_0^-\perp A_-)=819787,\qquad
 \vfour(\Qstar\perp\Hp)=787019.
\]
Thus this is a counterexample to blockwise positivity inferred from
maximum \(s\) alone, not to \cref{thm:classification}.
The negative lower-equality sign in \cref{lem:no-comp} cannot be omitted.
\end{example}

\begin{lemma}[The polar obstruction survives target changes]
\label{lem:mobius}
For a nondegenerate scalar pencil, a target basis change
\(\begin{psmallmatrix}a&b\\c&d\end{psmallmatrix}\in\GL_2(\F_2)\)
changes \(A=B_1^{-1}B_2\), up to input similarity, to
\begin{equation}\label{eq:mobius}
 (aI+bA)^{-1}(cI+dA).
\end{equation}
The operator is semisimple before this change if and only if it is
semisimple afterwards.
\end{lemma}
\begin{proof}
The new polar pair is \(B_1(aI+bA), B_1(cI+dA)\).
Its first member is nondegenerate by hypothesis, proving both the
formula and the invertibility of the denominator.
A semisimple operator diagonalizes over an algebraic closure, so any
defined rational function of it is semisimple. Applying the inverse
target transformation proves the converse. Input similarity also
preserves semisimplicity.
\end{proof}

\begin{example}[An odd-half-dimension lower equality without the core]
\label{ex:odd}
Let \(L=\F_8=\F_2[t]/(t^3+t+1)\), \(d=t^2+t+1\).
In \(V=L^2\oplus K^2\), with target \(W=\F_2^2\), take
\begin{align*}
 Q_{\mathrm{odd}}(x,y,u,v)=
 \bigl(&\Tr_8(xy)+\Tr(uv),\\
       &\Tr_8(txy+x+dy)+\Tr(\alpha uv)\bigr).
\end{align*}
By \cref{lem:F8,lem:H}, its scalar polars are nondegenerate,
its original sums are \(8\cdot4=32\), and its dual character sum is
\(-8\cdot16=-128\). Thus \(n=5\), the group order is \(4096\),
and
\[
 s=1120,\qquad \vfour=M_5-2^{10}=57467-1024=56443.
\]
It attains the odd lower endpoint in \cref{thm:bounds}.

Its \(A\) has a pair of simple \(t^3+t+1\)-blocks and a pair of
simple \(p\)-blocks, and is semisimple: under the trace pairing on
\(\F_8^2\), the two scalar polars correspond to multiplication by
\(1\) and \(t\), so \(A\) is multiplication by \(t\) on both
coordinates; on \(K^2\) it is multiplication by \(\alpha\).
The core \(\Qstar\) has a pair of \(p^2\)-blocks and is not semisimple.
If \(Q_{\mathrm{odd}}\) had a common orthogonal factor strictly
equivalent to \(\Qstar\), the corresponding restriction of its \(A\)
would be nonsemisimple after a common target change.
This contradicts \cref{lem:mobius}. It therefore has no such factor,
even with an unrestricted remaining summand.
The same polar obstruction applies to the odd equality types in
\cref{prop:equality}(a); it does not classify their complete square maps.
\end{example}

\section{Connections with existing results}\label{sec:context}
The quadratic-map and central-product correspondences used in
\cref{prop:strict,lem:central} are standard. For the latter the precise
source is \citep[Lemma 2.4, pp.~6--7]{KK12}, numbered according to its
arXiv v1 PDF. No strongly real hypothesis is required there.

The maximum-\(s\) branch is exactly the totally orthogonal branch
within \eqref{eq:class}. There are \(N\) linear characters with
Frobenius--Schur indicator \(+1\), and three nonlinear characters
\(\chi_\lambda\) of degree \(h\), inflated from the three extraspecial
quotients \citep[Proposition 3.5 and Remark 3.6, p.~8]{KK15}.
Their squared degrees sum to \(N+3h^2=|G|\), proving exhaustiveness.
The indicator definition yields
\[
 \nu(\chi_\lambda)
 =\frac1{4N}\sum_{g\in G}\chi_\lambda(g^2)
 =\frac{\Gauss(\lambda Q)}h.
\]
Thus all indicators are \(+1\) exactly when \(s=N+3h\).
This is the specialization of the Arf criterion restated in
\citep[Theorem 4.1, p.~12]{KK12}, where it is attributed to
Zahinda's Theorem 3.5. We use that attribution through the cited source,
not as a claim to have checked every result of the original paper.

For binary output dimension two, the three actual scalar duals are
distinct and nonzero: equality of two would contradict Fourier
inversion, and the dual of a bent function is bent. Their set, with
zero adjoined, is a two-dimensional function space if and only if their
sum \(D\) is identically zero. Consequently \(D\equiv0\) is exactly
vectorial dual-bent closure in \citep[Definition 1, p.~309]{CMP18} in
this situation. Any relabeling of the three nonzero components comes
from \(\GL_2(\F_2)\). The matrix identity
\[
 B_1^{-1}+B_2^{-1}=B_1^{-1}(B_1+B_2)B_2^{-1}
\]
makes the sum of two quadratic duals bent, but does not identify it
with the third dual.

On the maximum-\(s\) branch, \(D\equiv0\) supplies the same-label
vectorial dual and satisfies Condition A of \citep[p.~9]{WFW23}:
take \(p=2\), input dimension \(2n\), output dimension \(2\leq n\),
regularity coefficient \(\epsilon=1\), and \(F=Q\), with \(F(0)=0\).
For \(I=\{0\}\), \citep[Theorem 6, equation (35), pp.~23--24]{WFW23}
gives a partial difference set \(E=Q^{-1}(0)\setminus\{0\}\) with
\[
 z=\frac{N+3h}{4},\qquad k=z-1,\qquad
 \lambda=\frac N{16}+\frac{3h}{4}-2.
\]
If \(L\) counts the totally zero binary planes, the difference-set
parameter counts their six ordered bases, giving \(6L=k\lambda\).
Therefore the elementary lift count in \cref{lem:count} gives
\[
 \vfour=16L+6z-5=M_n+2^{3n-2}.
\]
The \(D\equiv0\) upper endpoint is thus a direct corollary. Binary
regularity does not assert a positive zero-frequency Gauss sum on
every orthogonal restriction. Neither \(C=0\) nor \(\Gauss(D)=0\)
can replace the condition \(D\equiv0\).

The lower-equality maps have \(\operatorname{rank}C=4\), so they are
not themselves vectorial dual-bent. The completed-class result of
\citep[Theorem 1, p.~311]{CMP18} concerns sets of scalar functions,
not a common strict equivalence of their simultaneous refinements.
Even allowing the full extended affine action
\(\widetilde Q(x)=TQ(Sx+a)+Lx+b\), the quadratic polars transform as
\(\widetilde B_\lambda=S^{\mathsf T}B_{T^*\lambda}S\), whence
\(\widetilde C=S^{-1}CS^{-\mathsf T}\). Translation and linear terms
do not change these polars. Thus rank four also prevents a simultaneous
extended-affine representative from becoming a closed rank-zero source.
This does not deny separate membership of scalar components in a
completed class. The group equivalence throughout remains
\eqref{eq:strict}, which is more restrictive.

The polar decompositions are instances of alternating-pencil theory.
\citep[Theorem 1.2]{LM18} enumerates isoclinism classes; its
Theorems 3.7--3.8 also impose exponent \(p\). We do not use them as
isomorphism classifications for the exponent-four groups here.
In characteristic two, homogeneous binary quadratic polynomials can
represent linear functions by diagonal squares. The local and
mixed-length framework in \citep[Section 4, pp.~15--18]{PFMR14}
therefore retains the full refinements. The equality module here is
\(R^2\oplus(R/(\eps))^{2r}\); its transformation formulas and the
classical Arf criterion underlie the short local normalizations.
Our explicit normalization does not depend on Proposition 7 of the
saved nineteen-page version of \citep{MPG13}. A counterexample to that
particular printed statement is given in \cref{app:MPG}; no conclusion
about the whole algorithm or other versions is inferred from it.

\appendix
\section{The common Maiorana--McFarland coordinates}\label{app:MM}
This appendix gives the parameter correspondence behind
\cref{lem:core-phase}. Let
\[
 E=K[\beta],\qquad \beta^2+\beta=\alpha,
 \qquad\Tr_{E/K}(\beta)=1.
\]
Thus \(E=\F_{16}\). Put
\[
 X=(b+d)+\beta d,\quad Y=a+\beta c,\qquad
 U=\upsilon+\beta\sigma,\quad
 V=(\tau+\rho)+\beta\rho.
\]
Direct multiplication gives
\[
 \Tr_{E/K}(XY)=ad+bc,\qquad
 \Tr_{E/\F_2}(UX+VY)
 =\Tr(\rho a+\sigma b+\tau c+\upsilon d).
\]
Write \(\mathcal H(Y)=\Tr(ac)+fa+ga^2+\ell c+mc^2\), interpreted
as a \(K\)-valued function of \(Y=a+\beta c\). For \(z\in K^\times\),
define
\begin{align*}
 h_z&=(z\theta+z^2j^2)+\beta(z\eta+z^2k^2),\quad
 \pi_z(Y)=zY+h_z,\qquad g_z(Y)=z\beta\mathcal H(Y).
\end{align*}
The trace identity \(\Tr(zw^2)=\Tr(z^2w)\) then gives exactly
\[
 \Tr(zQ^{(0)})
 =\Tr_{E/\F_2}\bigl(X\pi_z(Y)+g_z(Y)\bigr).
\]
Each \(\pi_z\) is an affine permutation of \(E\). Apply
\citep[equation (4), p.~313]{CMP18} with \(p=2\), \(m=4\), \(s=1\),
component parameter \(1\), and the trace pairing displayed above.
The actual phase is
\[
 \Tr_{E/\F_2}\bigl(V\pi_z^{-1}(U)
                     +g_z(\pi_z^{-1}(U))\bigr).
\]
Here \(\pi_z^{-1}(U)=a_z+\beta c_z\), with \(a_z,c_z\) from
\eqref{eq:az}. Substitution gives \eqref{eq:raw-core} including
\(E_z(0)\), for every one of the \(2^{16}\) refinements, not only the
\(96\) selected maps. When \(\eta=\theta=j=k=0\), the whole
\(K\)-valued core has the common representation
\(\Tr_{E/K}(XY+\beta\mathcal H(Y))\), with a generally nonzero
arbitrary function \(g\). The later \(g=0\) hypotheses of
\citep[Theorem 2]{CMP18} are not needed for equation (4).

\section{A version-specific vanishing-pencil example}\label{app:MPG}
We refer only to the nineteen-page ePrint revision of
\citep{MPG13}, Proposition 7 on p.~14. The consulted PDF has SHA-256
\texttt{6ce819d937f24089cf0e81f508129808}\allowbreak\texttt{e681346464a1c75cb5642981ef62aa1d}.
Its statement sends any vanishing
pencil to characteristics \((0,0,H^m,0)\), \(m\geq d/2\).
The following example checks the claim at the smallest nonzero
nilpotent parameter. It does not belong to the semi-extraspecial class
of \cref{sec:main} and plays no role in the proofs of the main results.

Take base and residue fields \(\F_2\), local length \(d=2\),
\(R=\F_2[H]/(H^2)\), \(t=1+H\), and
\(\lambda(a+Hb)=b\). In coordinates \(x=a+Hb\), \(y=c+Hd\), set
\[
 q_\infty=ad+bc,\qquad
 q_0=ad+bc+ac+a^2+c^2.
\]
Both polars are nondegenerate. In the block order \((b,a;d,c)\),
the multiplication matrix and trace matrix are
\[
 M=\begin{pmatrix}1&1\\0&1\end{pmatrix},\qquad
 T=\begin{pmatrix}0&1\\1&0\end{pmatrix},\qquad
 TM=\begin{pmatrix}0&1\\1&1\end{pmatrix}.
\]
The minimal and characteristic polynomials of \(M\) are both
\((X+1)^2\), so this is a cyclic regular pencil of the required
primary type. The first form has no diagonal terms, and the second
has the two diagonal characteristics \(H,H\): the full tuple is
\((0,0,H,H)\). Moreover the invertible binary change
\[
 (a,b,c,d)\longmapsto(a,b+c,c,d+a+c)
\]
interchanges the two quadratic forms. Indeed its pullback sends
\(q_\infty\) to \(q_0\), and adding \(ac+a^2+c^2\) sends the latter
back to \(q_\infty\). Thus both ordered pencils are semi-vanishing,
which is exactly the printed vanishing condition.

For the input pair, simultaneous vanishing requires
\(ac+a+c=0\), whose only solution in \(\F_2^2\) is \(a=c=0\).
There are therefore four common zeros. The printed targets at \(d=2\)
have second form either
\(ad+bc+ac+a^2\) or \(ad+bc+ac\). In the first case the permitted
\((a,c)\) are \((0,0),(0,1),(1,1)\); in the second they are
\((0,0),(0,1),(1,0)\). Each gives respectively four, two, and two
choices of \((b,d)\), so each target has eight common zeros.
An invertible input change preserves this number, as does even a
common invertible target change. Hence the printed reduction is false
for this example over the original field. No field extension or
affine operation is part of that fixed-field isomorphism statement.
This observation is restricted to that version and proposition; it
does not invalidate its other local transformation formulas or claim
anything about originality of the main results here.

\section*{Use of generative AI}
Generative AI tools, including OpenAI Codex/ChatGPT and the system
recorded as Muse Spark 1.3, contributed to mathematical exploration
and proof development, computational checks, literature review, and
manuscript drafting. This included substantive mathematical steps and
initial text generation, not only language editing. Exact historical
deployment versions are not fully documented. Internal automated checks
are not external peer review.

\begingroup
\raggedright
\begin{thebibliography}{8}
\bibitem[Abdukhalikov et~al.(2024)]{AFH24}
Abdukhalikov, K., Feng, R., Ho, D. (2024).
Quadratic bent functions and their duals.
\emph{Applicable Algebra in Engineering, Communication and Computing}
\textbf{35}(4), 507--523.
\url{https://doi.org/10.1007/s00200-022-00564-5}.

\bibitem[\c{C}e\c{s}melio\u{g}lu et~al.(2018)]{CMP18}
\c{C}e\c{s}melio\u{g}lu, A., Meidl, W., Pott, A. (2018).
Vectorial bent functions and their duals.
\emph{Linear Algebra and its Applications} \textbf{548}, 305--320.
\url{https://doi.org/10.1016/j.laa.2018.03.016}.

\bibitem[Kaur and Kulshrestha(2012)]{KK12}
Kaur, D., Kulshrestha, A. (2012). Strongly real special 2-groups.
arXiv:1210.3790v1.
\url{https://arxiv.org/abs/1210.3790v1}.
Cited numbering and pages refer to this version.

\bibitem[Kaur and Kulshrestha(2015)]{KK15}
Kaur, D., Kulshrestha, A. (2015). Characters of real special 2-groups.
arXiv:1510.06583v1.
\url{https://arxiv.org/abs/1510.06583v1}.

\bibitem[Lewis and Maglione(2020)]{LM18}
Lewis, M. L., Maglione, J. (2020).
Enumerating isoclinism classes of semi-extraspecial groups.
\emph{Proceedings of the Edinburgh Mathematical Society}
\textbf{63}, 426--442.
\url{https://doi.org/10.1017/S0013091519000531}.
Cited numbering and PDF pages refer to arXiv:1806.10511v2 (2018),
\url{https://arxiv.org/abs/1806.10511v2}.

\bibitem[Macario-Rat et~al.(2013)]{MPG13}
Macario-Rat, G., Pl\^ut, J., Gilbert, H. (2013).
New Insight into the Isomorphism of Polynomial Problem IP1S and Its
Use in Cryptography. IACR Cryptology ePrint Archive, Report 2013/799.
\url{https://eprint.iacr.org/2013/799.pdf}.
Cited statements refer to the nineteen-page revision,
not the seventeen-page conference archive.

\bibitem[Pl\^ut et~al.(2014)]{PFMR14}
Pl\^ut, J., Fouque, P.-A., Macario-Rat, G. (2014).
Solving the ``Isomorphism of Polynomials with Two Secrets'' Problem
for All Pairs of Quadratic Forms. arXiv:1406.3163v3.
\url{https://arxiv.org/abs/1406.3163v3}.

\bibitem[Wang et~al.(2023)]{WFW23}
Wang, J., Fu, F.-W., Wei, Y. (2023).
Bent Partitions, Vectorial Dual-Bent Functions and Partial Difference
Sets. arXiv:2301.00581v1.
\url{https://arxiv.org/abs/2301.00581v1}.
\end{thebibliography}
\endgroup
\end{document}
